\section{Related Work}
\label{sec:related}

\para{Logical fallacy detection.}
The task of automated logical fallacy detection was formalized by \citet{jin2022logical}, who introduced the \logicdataset dataset with 13 fallacy categories drawn from educational sources. Subsequent benchmarks expanded coverage: \citet{hong2024closer} created the \fallacies dataset with 232 types in a hierarchical taxonomy, while \citet{helwe2024mafalda} consolidated multiple sources into the \mafalda benchmark with span-level annotations across 25 types. \citet{pan2024llms} evaluated LLMs as zero-shot fallacy classifiers, finding that multi-round prompting with chain-of-thought reasoning improves detection. \citet{yeh2024cocolofa} contributed ecologically valid examples from news comments. Across these studies, GPT-4 consistently achieves the highest detection rates (87.7\% on binary identification per \citet{hong2024closer}), though fine-tuned BERT models remain competitive for in-domain settings.

\para{LLM self-evaluation and reasoning errors.}
A growing body of work examines whether LLMs can identify errors in their own reasoning. \citet{tyen2024llms} decomposed self-correction into mistake finding and output correction, demonstrating that LLMs fail at the former (GPT-4 at 52.87\% accuracy) but succeed at the latter given oracle error locations. \citet{liu2024self} documented self-contradictory reasoning rates up to 50\% and showed that GPT-4 detects these at only 52.2\% F1, substantially below human performance (66.7\% F1). \citet{payandeh2023susceptible} identified a self-detection paradox where GPT-3.5 generates fallacies it cannot detect, and found that GPT-4 is 69\% more susceptible to fallacious arguments than logical ones in debate settings. These results collectively suggest that self-evaluation is a bottleneck, but none of these studies controlled for whether \emph{attribution framing} (labeling reasoning as ``yours'' versus ``another model's'') biases detection independently of content.

\para{Chain-of-thought faithfulness.}
The reliability of chain-of-thought reasoning as a window into model computation is itself contested. \citet{lyu2023faithful} showed that translating reasoning into symbolic language executed by deterministic solvers outperforms standard CoT on 8/10 benchmarks, suggesting standard CoT is unfaithful to actual reasoning. \citet{arcuschin2025measuring} found that faithfulness does not correlate with plausibility ($r = 0.15$), implying that plausible-sounding reasoning may contain hidden errors. These findings motivate our study: if CoT reasoning can be unfaithful, models may produce and fail to detect fallacies within their own chains of reasoning.

\para{Scaling and human-like fallacy patterns.}
\citet{richardson2025theory} tested 38 models on 383 reasoning problems and found that more capable models produce \emph{more human-like} fallacies ($\rho = 0.360$, $p = 0.027$), while overall correctness shows no correlation with capability ($r = 0.004$). \citet{zhai2025ruozhibench} found that even the best model achieves only 62\% on deceptive reasoning tasks versus a human upper bound above 90\%. These results suggest that scaling alone does not eliminate systematic blindspots. Our work complements these findings by asking whether the source of self-evaluation failure lies in attribution bias or in absolute detection difficulty.

\para{Positioning of our work.}
Unlike prior studies that evaluate fallacy detection in a single condition, we systematically vary source attribution while holding content constant. Unlike \citet{tyen2024llms}, who studied step-level error finding in mathematical reasoning, we focus on logical fallacy detection across 13 categories. Unlike \citet{hong2024closer}, who documented category-specific variation without examining attribution effects, we directly test whether labeling reasoning as ``self-generated'' changes detection rates. Our cross-model evaluation design extends beyond attribution framing to test whether genuine self-evaluation differs from cross-model evaluation in practice.
